%------------------------
% Resume Template
% Author : Anubhav Singh
% Github : https://github.com/xprilion
% License : MIT
%------------------------

\documentclass[a4paper,20pt]{article}

\usepackage{latexsym}
\usepackage[empty]{fullpage}
\usepackage{titlesec}
\usepackage{marvosym}
\usepackage[usenames,dvipsnames]{color}
\usepackage{verbatim}
\usepackage{enumitem}
\usepackage{hyperref}
\usepackage{fancyhdr}
\usepackage[UTF8, fontset=fandol]{ctex}

\pagestyle{fancy}
\fancyhf{} % clear all header and footer fields
\fancyfoot{}
\renewcommand{\headrulewidth}{0pt}
\renewcommand{\footrulewidth}{0pt}

% Adjust margins
\addtolength{\oddsidemargin}{-0.530in}
\addtolength{\evensidemargin}{-0.375in}
\addtolength{\textwidth}{1in}
\addtolength{\topmargin}{-.45in}
\addtolength{\textheight}{1in}

\urlstyle{rm}

\raggedbottom
\raggedright
\setlength{\tabcolsep}{0in}

% Sections formatting
\titleformat{\section}{
  \vspace{-10pt}\scshape\raggedright\large
}{}{0em}{}[\color{black}\titlerule \vspace{-6pt}]

%-------------------------
% Custom commands
\newcommand{\resumeItem}[2]{
  \item\small{
    \textbf{#1}{: #2 \vspace{-2pt}}
  }
}

\newcommand{\resumeItemWithoutTitle}[1]{
  \item\small{
    {\vspace{-2pt}}
  }
}

\newcommand{\resumeSubheading}[4]{
  \vspace{-1pt}\item
    \begin{tabular*}{0.97\textwidth}{l@{\extracolsep{\fill}}r}
      \textbf{#1} & #2 \\
      \textit{#3} & \textit{#4} \\
    \end{tabular*}\vspace{-5pt}
}

\newcommand{\resumeSubItem}[2]{\resumeItem{#1}{#2}\vspace{-3pt}}

\newcommand{\resumeProjectItem}[3]{
  \item
    \begin{tabular*}{0.97\textwidth}{l@{\extracolsep{\fill}}r}
      \textbf{#1} & \textit{#2} \\
    \end{tabular*}
    \vspace{-2pt}
      #3
    \vspace{-2pt}
}

\renewcommand{\labelitemii}{$\circ$}

\newcommand{\resumeSubHeadingListStart}{\begin{itemize}[leftmargin=*]}
\newcommand{\resumeSubHeadingListEnd}{\end{itemize}}
\newcommand{\resumeItemListStart}{\begin{itemize}}
\newcommand{\resumeItemListEnd}{\end{itemize}\vspace{-5pt}}

%-----------------------------
%%%%%%  CV STARTS HERE  %%%%%%

\begin{document}

%----------HEADING-----------------
\begin{tabular*}{\textwidth}{l@{\extracolsep{\fill}}r}
   \textbf{{\LARGE 聂家明}} \\
   \href{https://www.linkedin.com/in/jmnie/}{LinkedIn: in/jmnie}  & 邮箱:\href{mailto:jiaming.nie13@gmail.com}{jiaming.nie13@gmail.com} \\
  & 电话: (+86)-135-2171-3097
\end{tabular*}

%-----------EDUCATION-----------------
\section{教育背景}
  \resumeSubHeadingListStart
    \resumeSubheading
      {伍斯特理工学院}{马萨诸塞，美国}
      {硕士 - 计算机科学}{2017.8 - 2019.12}
    \resumeSubheading
      {利物浦大学}{利物浦，英国}
      {学士 - 电气工程}{2015.8 - 2017.6}
    \resumeSubheading
      {西交利物浦大学}{苏州，中国}
      {学士 - 电气工程}{2013.8 - 2015.8}
    \resumeSubHeadingListEnd

\vspace{2pt}
\section{技能}
	\resumeSubHeadingListStart
	\resumeSubItem{语言}{~~~Python, JavaScript/TypeScript, Java}
    \resumeSubItem{AI/LLM}{~~~LLM Fine-Tuning (SFT/RFT), RAG}
    \resumeSubItem{后端/部署}{~~~FastAPI, Flask, Node.js, Redis, MongoDB, Docker, AWS, Prometheus, Grafana}
	\resumeSubHeadingListEnd

\vspace{4pt}
\section{工作经历}
  \resumeSubHeadingListStart
  \vspace{4pt}
  \resumeSubheading{Youlify}{远程}
    {创始软件工程师}{2025.1 - 至今}
    \resumeItemListStart
        \item{作为医疗 AI 初创公司 \textbf{Founding Engineer}，围绕诊所理赔自动化、拒赔审查和 \textbf{EHR} 数据接入等业务场景，负责核心 \textbf{GenAI} 能力与后端基础设施建设，推动 AI 项目从原型验证走向生产落地。}
        \item{独立设计并实现医疗计费代码抽取的端到端 \textbf{LLM} 微调与评测流水线，基于 2,000+ 份临床记录构建 \textbf{SFT/RFT} 数据集，自动化完成数据处理、训练任务编排、检查点跟踪、效果评测与成本分析。}
        \item{建设拒赔审查与申诉信生成的后端工作流，将理赔单、\textbf{EOB}、保险数据等转化为可追溯上下文，完成模型路由、异步任务触发、状态回传和业务状态管理，支撑长链路理赔自动化任务。}
        \item{主导单体后端向 \textbf{FastAPI} 微服务架构迁移，拆分诊所管理、Portal API、鉴权和 \textbf{GenAI} 服务，并搭建 \textbf{Prometheus + Grafana} 可观测体系，提升系统稳定性与迭代效率。}
        \item{引入 \textbf{Apache Airflow} 编排患者 \textbf{EHR} 数据与临床文档同步长任务，打通 Portal、Integration 与 Automation Jobs 的任务触发、取消、状态持久化与恢复链路，支撑多诊所、多 \textbf{EHR} 系统的数据同步与进度追踪。}
    \resumeItemListEnd

    \vspace{4pt}
  \resumeSubheading{北京峥研软件有限责任公司（peerup）}{上海, 中国}
    {软件开发工程师}{2023.6 - 2025.1}
    \resumeItemListStart
        \item{使用 \textbf{FastAPI} 为笔记产品中的 \textbf{AI} 助手和聊天机器人构建后端服务，支持流式响应、文件理解、编辑器内动作触发等核心能力。}
        \item{构建基于 \textbf{Milvus} 与 \textbf{Elasticsearch} 的混合检索 \textbf{RAG} 流水线，整合公共知识库、用户内容与外部数据源，为回答生成提供可追溯上下文。}
        \item{在 \textbf{FastAPI} 生态中使用 \textbf{LangChain} 与 \textbf{LlamaIndex} 封装检索、上下文组装和外部服务调用能力，支撑多步骤 \textbf{AI} 功能在后端稳定执行。}
        \item{维护 \textbf{Docker} 容器化与自动化 \textbf{CI/CD} 流水线，支撑 AI 功能稳定发布与持续迭代。}
      \resumeItemListEnd

    \vspace{4pt}
    \resumeSubheading{Oqton}{上海, 中国}
    {软件开发工程师}{2021.10 - 2023.2}
    \resumeItemListStart
        \item{面向工业制造场景开发 \textbf{Python} 服务，与工业设备进行交互，采集机器数据并通过 \textbf{API} 下发控制指令，参与企业级系统集成、现场问题定位与交付支持。}
        \item{在 \textbf{Linux} 网关上部署并维护 \textbf{Docker} 应用，协助排查跨 \textbf{Elasticsearch} 与 \textbf{Kafka} 系统的生产环境问题。}
        \item{使用 \textbf{Node.js}、\textbf{Express} 和 \textbf{React} 开发机器配置与数据管理相关功能，支持工业软件业务流程落地。}
      \resumeItemListEnd

    \vspace{4pt}
    \resumeSubheading
		{微福思软件科技有限公司(Micro Focus)}{上海, 中国}
		{软件开发工程师}{2020.7 - 2021.9}
		\resumeItemListStart
              \item{使用 \textbf{JavaScript}、\textbf{Node.js} 与 \textbf{C++} 开发 Web 组件检测相关自动化测试能力。}
              \item{实现跨域组件检测功能，并优化内部案例处理工作流。}
		\resumeItemListEnd

\resumeSubHeadingListEnd

%-----------PROJECTS-----------------
\vspace{-5pt}
\section{代表项目}
\resumeSubHeadingListStart

\vspace{4pt}
\resumeProjectItem{医疗计费 LLM 微调工作流}{2026.2 - 至今}{
\resumeItemListStart
    \item{设计并实现端到端 \textbf{OpenAI} 托管微调流水线，处理 2,000+ 份临床记录，构建 \textbf{SFT} 与 \textbf{RFT} 数据集，用于医疗计费代码抽取。}
    \item{通过 \textbf{OpenAI Fine-Tuning API} 编排 \textbf{GPT-4.1 SFT} 和 \textbf{o4-mini RFT} 训练任务，实现文件上传、状态轮询、检查点跟踪和实验元数据管理自动化。}
    \item{为 \textbf{RFT} 自研 \textbf{Python} grader，结合加权编辑距离与领域惩罚项（\textbf{CPT} 前缀、\textbf{ICD-10}、modifier）优化结构化 \textbf{JSON} 输出质量。}
    \item{建立完整评测框架，追踪宏观/微观指标，并结合 token 成本估算与 checkpoint 分析对比基线模型效果。}
\resumeItemListEnd
}

\vspace{4pt}
\resumeProjectItem{医疗 SaaS 数据同步与实时任务平台}{2025.1 - 至今}{
\resumeItemListStart
    \item{设计并实现基于 \textbf{Apache Airflow} 的患者 \textbf{EHR} 数据同步编排，将 Athena、Practice Fusion 等传统医疗系统的长耗时拉取任务迁移为可触发、可取消、可追踪的 \textbf{DAG}。}
    \item{使用 \textbf{FastAPI + Gunicorn/Uvicorn} 构建高并发后端 API，并通过 \textbf{WebSocket} 向前端推送抓取进度、错误状态和任务结果。}
    \item{基于 \textbf{Redis} 分布式锁、任务队列状态和 \textbf{MongoDB} 多租户存储设计幂等处理流程，保障诊所租户间的数据隔离与失败重试。}
    \item{配合客户接入和上线排障，沉淀可观测日志、任务状态查询和异常恢复流程，缩短新诊所接入与数据验证周期。}
\resumeItemListEnd
}

\vspace{4pt}
\resumeProjectItem{AI 助手与知识增强系统}{2023.9 - 2025.1}{
\resumeItemListStart
    \item{主导开发基于 \textbf{FastAPI} 的 AI 助手后端，支持 \textbf{SSE} 流式响应、复杂指令执行、文件上传问答与编辑器内动作触发。}
    \item{构建基于 \textbf{Milvus} 向量库和 \textbf{Elasticsearch} 的混合检索系统，整合公共知识库、用户私有知识库与第三方网络结果，增强回答准确性和上下文相关性。}
    \item{使用 \textbf{LangChain} 与 \textbf{LlamaIndex} 封装检索、上下文构造和外部服务调用能力，支撑后端多步骤 \textbf{AI} 任务执行。}
    \item{设计 \textbf{Redis} 会话管理与权限隔离机制，支持消息、对话、附件映射缓存及基于用户层级的数据检索控制。}
\resumeItemListEnd
}

\resumeSubHeadingListEnd

\end{document}
